% math-tagging-fixture.tex
% Focused PDF/UA-1 fixture for checking Formula tagging and math /Alt behavior.
% Compile with LuaLaTeX. PAC/Acrobat validation still must be run separately.
\DocumentMetadata{
  tagging=on,
  pdfstandard=ua-1,
  lang=en-US,
  tagging-setup={
    math/alt/use
  }
}

\documentclass{../ndsudisq}

\disquisitiontitle{Math Tagging Fixture}
\disquisitionauthor{Template Test}
\documenttype{Dissertation}
\degree{Doctor of Philosophy}
\majordepartmentlabel{Major Department}
\program{Test Program}
\degreeoption{}
\finaldefense{May 2026}
\approvaldate{May 18, 2026}
\departmentchair{Department Chair}
\committeechair{Committee Chair}
\committeeone{Committee One}
\committeetwo{Committee Two}
\committeethree{Committee Three}
\subjectinfo{Formula tagging fixture}
\keywordsinfo{NDSU, disquisition, Formula tags, math alternative text}
\disquisitionlanguage{en-US}
\includelistofequations
\numberwithin{equation}{chapter}

\begin{document}
\makefrontmatter
\mainmattersetup

\chapter{Math Tagging Fixture}
\label{chap:math-tagging-fixture}

% These examples are present to test Formula tagging and math /Alt behavior
% under the PDF/UA-1 validation target. They should remain LaTeX math rather
% than screenshots or image replacements.
This paragraph contains the inline formula \(S=\sum_{i=1}^{n} w_i x_i\).
It also contains the inline variables \(\bar{x}\), \(x_i\), and \(n\).

% The \mathalt command preserves the template's existing display-equation
% description workflow for numbered display equations while math/alt/use
% supplies a fallback for Formula tags.
\mathalt{Sample mean equals one over n times the sum from i equals one to n of x sub i.}
\begin{equation}
\bar{x}=\frac{1}{n}\sum_{i=1}^{n}x_i
\end{equation}

% Include an amsmath environment so the fixture covers display math beyond the
% base equation environment.
\begin{align}
S &= \sum_{i=1}^{n} w_i x_i \\
\bar{x} &= \frac{1}{n}\sum_{i=1}^{n}x_i
\end{align}

\end{document}
